% ==================== PART 2 — System design (human-written) ====================
% SKELETON for Andy, ~1 page. Part 2 = describe the AGENT / harness / tools / research loop + commentary on significance.
% It is the META layer (how the research was done with the agent), NOT a re-telling of the Part 1 research.
% <=4pp is a MAX; aim for ~1 page. Fill bullets in your own voice. [BRACKETS] = things only you can fill in.
% A fuller prose version of the verification/exploit story is in part2_prose_reference.tex if you want to expand.

\section{The agent and how we used it}
\label{sec:agent}
% Two distinct roles for "the agent" -- name them.
\begin{itemize}\itemsep2pt
  \item \textbf{The evolution loop (the main loop).} An island program-evolution framework (ShinkaEvolve~\citep{shinkaevolve})
  with two Claude models as the \emph{proposal} agents that wrote and mutated candidate tester programs:
  \texttt{claude-sonnet-5} (breadth) and \texttt{claude-opus-5} (depth). Two islands; a fixed API-cost cap ([\$X]); [N]
  generations. This loop did the research: it evolved the testers, ran the experiments, and drafted Part~1.
  \item \textbf{The interactive agent.} I drove Claude Code ([model / effort], e.g.\ Opus~4.8) to steer the project:
  reality-checking numbers, running an adversarial review of my own claims, and drafting these Parts. [Say roughly how
  many times you checked in / the main decision points you steered: e.g.\ choosing the evaluator, the reframe, the final
  submission.]
  \item \textbf{Division of labor.} The loop does the research; I designed the defect oracles and the protocol, decided
  what to measure, verified every number against regenerable logs, and made the interventions below.
\end{itemize}

\section{Harness and tools}
\label{sec:harness}
\begin{itemize}\itemsep2pt
  \item \textbf{Substrates:} the genuine published simulators, all pure-JAX (Kinetix/jax2d, MiniGrid, MinAtar, Jumanji
  Sokoban, Gymnax, Craftax). PPO baselines run in the same harness on the same feature vector.
  \item \textbf{Fitness / compute:} candidates scored by the fraction of witness states reached at a fixed budget, on a
  multi-engine search split (seeds $\{0,1,2\}$), run on cluster GPUs (Hyak/SLURM); orchestration on a laptop.
  \item \textbf{Prompt \& feedback:} the prompt gives the four-hook interface, the scoring formula, and hard constraints
  (pure JAX, fixed shapes, in-range actions, ``runs in a few seconds per level''); each candidate gets back its scalar
  fitness plus its worst targets and a paired $95\%$ CI against a frozen reference.
  \item \textbf{Verifier:} replay-from-reset of every counted hit, a frozen hashed protocol, sealed level and defect-class
  splits, and a reachability calibration (details in the research-loop step below).
\end{itemize}

\section{The research loop, and the one check that mattered}
\label{sec:loop}
\begin{itemize}\itemsep2pt
  \item \textbf{Cadence:} we froze the oracles and a hashed protocol, pointed the evolution loop at a multi-engine
  fitness, reality-checked its numbers (replay + exclusions), ran an adversarial review, reframed the claim to what the
  data supports, and wrote the paper. [Fill in your check-in count / where you intervened.]
  \item \textbf{Verification is load-bearing, not decorative.} The loop's first strong ``discovery'' inflated its own
  fitness by \textbf{+0.11} through a concurrency race in the harness (physically impossible states, nothing wrong in the
  candidate's code). It passed the fitness, both held-out splits, four design comparisons, and an LLM code-audit; only
  \textbf{replaying every hit from a clean reset} caught it. We rebuilt the harness, withdrew the tainted number, and
  corrected our own $N{=}3$ over-reach to $N{=}1$. [Expand this in your own words if you want; it is the strongest story.]
\end{itemize}

\section{Significance of the main result}
\label{sec:significance}
\begin{itemize}\itemsep2pt
  \item \textbf{Cost.} A written tester is a program that runs \emph{training-free}: evolve one transferable tester once
  and reuse it across games, versus training a fresh RL explorer for each. At a tester's real budget the written program
  wins, and RL only catches up given far more compute, so for the same testing outcome the written route is much cheaper
  to deploy.
  \item \textbf{Accessibility.} This lowers the bar for who can test a game. A small studio or an individual without an
  RL-training pipeline, but with access to an LLM, can run a written tester to explore their game for bugs, instead of
  standing up per-game reinforcement learning.
  \item \textbf{For autonomous research broadly.} Execution-grounded replay is what let us trust the loop's claims; it
  should live \emph{inside} the loop's fitness so such exploits are worthless by construction. That is the workshop's
  bottleneck (evaluating AI-generated results) made concrete.
\end{itemize}
